% !TEX root = ../bb_AiRa_Kappaleet.tex

% ö = \%c3\%b6
% ä = \%C3\%A4

\xdef\sectiontitle{\Isectiontitle} %aiheuttama
\TOCrefs

%%%%%%%%%%%%%%%
\begin{frame}[label=1_9_a]%[label=1_9_TOC1]
\frametitle{\texorpdfstring{\culine[\sectiontitleulinecolor]{\sectiontitle}}{\sectiontitle} }
\label{\TOCname}

\vspace{-0.5cm}

\begin{itemize}[leftmargin=0.5cm,itemsep=2mm]
    \item[1.1]<1-> \hyperlink{1_1_a}{\IaSubsectiontitle}
    \item[1.2]<1-> \hyperlink{1_2_a}{\IbSubsectiontitle}	
    \item[1.3]<1-> \hyperlink{1_3_a}{\IcSubsectiontitle}	
    \item[1.4]<1-> \hyperlink{1_4_a}{\IdSubsectiontitle}
    \item[1.5]<1-> \hyperlink{1_5_a}{\IeSubsectiontitle}
    \item[1.6]<1-> \hyperlink{1_6_a}{\IfSubsectiontitle}
    \item[1.7]<1-> \hyperlink{1_7_a}{\IgSubsectiontitle}
\end{itemize}

\vspace{-1cm}
\begin{itemize}[leftmargin=8.5cm,itemsep=2mm]
    \item[1.8]<1-> \hyperlink{1_8_a}{\IhSubsectiontitle}
	\item[1.9]<1-> \hyperlink{1_9_a}{\CDAlert[blue]{\IiSubsectiontitle}}
	\item[1.10]<1-> \hyperlink{1_10_a}{\IjSubsectiontitle}
	\item[1.11]<1-> \hyperlink{1_11_a}{\IkSubsectiontitle}
\end{itemize}

%\endofslide 
\end{frame}
%%%%%%%%%%%%%%%


\xdef\sectiontitle{\IiSubsectiontitle}

%%%%%%%%%%%%%%%
\begin{frame}%[label=1_9_a]
\frametitle{\texorpdfstring{\culine[\sectiontitleulinecolor]{\sectiontitle}}{\sectiontitle} }
\framesubtitle{Ominaisuudet ja peruskomponentit}

\appear{}
%\vspace{1.5mm}
\begin{minipage}{8.1cm}
\begin{itemize}[itemsep=1.7mm]
	\item Laserit ovat loistava esimerkki fotonikaasusysteemistä
	\item Ne ovat laitteita joissa suuren energiatiheyden omaavaa SMG-säteilyä \appear{on laserin kaviteetin sisällä}\appear{, josta tulee ulos \Alert[red]{on koherenttia ja yksisuuntaista monokromaattista valoa} }
\appear{\xdef\loc{east}%
\begin{tikzpicture}[remember picture,overlay] % anchor=south
    \node (A) [yshift=0cm,xshift=\figXshift, anchor=\loc] at (current page.\loc) {\includegraphics[page=1,height=8cm]{Figures_booklet/SOM_9_Statistical_mechanics_figures/SOM_Figure_21.png}}; %scale=\figscale. % 8cm max hei
\end{tikzpicture}}%

\item Laserin peruskomponentit ovat:
\begin{itemize}
\small
	\item \CDAlert[fuchsia]{kolmitasosysteemi}: \appear{jossa ylemmällä "2"~tasolla olevat atomit} \appear{ siirtyvät alemmalle tilalle  "1"~}\appear{prosessissa jonka stimuloi energiaeroa vastaavalla taajuudella värähtelevä fotoni}
\item \CDAlert[blue]{Optinen kaviteetti}: \appear{joka vahvistaa 'oikeanlaisia' fotoneita}
\item \CDAlert[red]{Pumppausmekanismi}: joka nostaa atomeita tilalle  "2"~\appear{jatkuvatoimisuuden aikaansaamiseksi}
\end{itemize}
\end{itemize}
\end{minipage}



\endofslide 
%\endofsection
\end{frame}
%%%%%%%%%%%%%%%

%%%%%%%%%%%%%%%
\begin{frame}
\frametitle{\texorpdfstring{\culine[\sectiontitleulinecolor]{\sectiontitle}}{\sectiontitle} }
\framesubtitle{Toimintaperiaatteet}
\appear{}

\begin{itemize}
	\item Laserin toimintaperiaatteet ovat käytännössä:

\item[] \begin{itemize}[itemsep=2.3mm] \small
	\item Väliaineen reunoilla on samansuuntaiset peilit
  \item Optisen akselin suhteen eri suuntaan etenevät fotonit \appear{eivät juurikaan stimuloi emissiota}
  \item \Alert[blue]{Akselin suuntaiset fotonit heijastelevat peilien välissä}\appear{, mikä mahdollistaa stimuloidun emission}
  \item Pituus $L$ \appear{määrittää seisovan aallon:} \appear{$L=n \lambda / 2$}
    % 6.5 cm = midway, 10 cm = 3/4 
\vspace{2.3mm}
\end{itemize}
\end{itemize}
\begin{minipage}{6.25cm}
\begin{itemize}
\item[] \begin{itemize}[itemsep=2.3mm]  \small
  \item Kaviteettiin muodostuu resonanssi
  \item \CDAlert[blue!90!black]{Akselin suuntainen valo vahvistuu}
  \item Erisuuntainen valo ei vahvistu
  \item \Alert[red]{Laserin aallonpituus kapenee suhteessa atomien spontaanin transition viivanleveyteen}
  \item Energiaa tuodaan kaviteettiin esim. 'optisella \Alert[fuchsia]{pumppauksella}'
\end{itemize}
\end{itemize}
  \end{minipage}

\xdef\loc{south east}
\begin{tikzpicture}[remember picture,overlay] % anchor=south
    \node (A) [yshift=-0.25*\figYshift,xshift=\figXshift, anchor=\loc] at (current page.\loc) {\includegraphics[page=1,width=6.7cm]{Figures_booklet/SOM_9_Statistical_mechanics_figures/SOM_Figure_23.png}}; %scale=\figscale. % 8cm max hei
\end{tikzpicture}


\endofslide 
%\endofsection
\end{frame}
%%%%%%%%%%%%%%%

%%%%%%%%%%%%%%%
\begin{frame}
\frametitle{\texorpdfstring{\culine[\sectiontitleulinecolor]{\sectiontitle}}{\sectiontitle} }
\framesubtitle{Sironta- ja absorptioprosessit}
\appear{}

\begin{minipage}{7cm}
\begin{itemize}[itemsep=2.7mm]
	% 6.5 cm = midway, 10 cm = 3/4 
\vspace{2mm}
\small
	\item \CDAlert[red]{Spontaania emissiota} \appear{tapahtuu kun korkeammalla tilalla oleva elektroni} \appear{tipahtaa alemmalle tilalle lyhyen ajan jälkeen (tilastollinen prosessi)}\appear{, emittoiden fotonin}

\item \CDAlert[blue]{Absorptiota} \appear{tapahtuu kun elektroni viritetään kor\-keammalle energiatilalle} \appear{ (systeemi absorboi fotonin)}

\item \CDAlert[fuchsia]{Stimuloidussa emissiossa}\appear{, korkeammalla tilalla oleva elektroni siirtyy läsnäolevan fotonin vaikutuksesta alemmalle tilalle.}\appear{ Fotonin energia vastaa täsmälleen elektronitilojen energiaeroa} 
 \item[] \begin{itemize}
	\item[=>] Syntyy toinen fotoni\appear{, jonka energia} \appear{(eli aallonpituus) }\appear{on identtinen läsnäolevan fotonin energian kanssa}\appear{, fotonien muutkin ominaisuudet ovat identtisiä}\appear{/korreloivat} \appear{keskenään}\appear{, kuten suunta}\appear{, \CDAlert[fuchsia]{vaihe}} \appear{ja polarisaatio} 
\end{itemize}
 
\end{itemize}
\end{minipage}

\xdef\loc{east}
\begin{tikzpicture}[remember picture,overlay] % anchor=south
    \node (A) [yshift=0cm,xshift=\figXshift, anchor=\loc] at (current page.\loc) {\includegraphics[page=1,height=8cm,trim={0 0 13cm 0},clip]{Figures_booklet/SOM_9_Statistical_mechanics_figures/SOM_Figure_22.png}}; %scale=\figscale. % 8cm max hei
\end{tikzpicture}


\endofslide 
%\endofsection
\end{frame}
%%%%%%%%%%%%%%%

%%%%%%%%%%%%%%%
\begin{frame}
\frametitle{\texorpdfstring{\culine[\sectiontitleulinecolor]{\sectiontitle}}{\sectiontitle} }
\framesubtitle{Sironta- ja absorptioprosessit}
\appear{}

\begin{minipage}{7.2cm}
\begin{itemize}\small 
	% 6.5 cm = midway, 10 cm = 3/4 
%\vspace{2mm}
	\item Kuinka kuvata prosesseja matemaattisesti?
\appear{Yksinkertaisuuden vuoksi tarkastellaan kaasumaista vahvistinainetta}\appear{ (energiatilatarkastelu yksinkertaisempaa \appear{kuin kiinteässä aineessa)}} 

\item Spontaanissa emissiossa\appear{ \CDAlert[red]{transitionopeuden} \CDAlert[red]{$R_{\text{spon}}$}} \appear{ täytyy olla verrannollinen \Alert[red]{virittyneiden ato\-mi\-en lukumäärään, $N_{2}$}}\appear{, jolloin verrannollisuuskertoi\-meksi saadaan $A_{\text{spon}}$} \appear{(}\Alert[red]{$\appear{R_{\text{spon}}} \equal \appear{A_{\text {spon}}} \appear{N_2}$}\appear{)}

\item \appear{\CDAlert[blue]{Absorptionopeus}} on verrannollinen alemmalla tasolla olevien hiukkasten lukumäärään \appear{$N_{1}$}\appear{, sekä oikean energian omaavien fotonien lukumäärään} \appear{$Y(\Delta E)$}\appear{, missä $\Delta E=E_{2}-E_{1}$.} \appear{\Alert[blue]{Verrannollisuusvakio on $B_{abs}$}} \appear{(}\Alert[blue]{$\appear{R_{abs}} \equal \appear{B_{abs}} \appear{N_{1}} \appear{Y(\Delta E)}$}\appear{)}

\item \CDAlert[fuchsia]{Stimuloitu emissio}\appear{ on verrannollinen ylemmän tilan omaavien vapaiden atomien lukumäärään}\appear{, sekä oikean energian omaavien fotonien lukumäärään}\appear{, jotka voivat ajaa emissioprossia} \appear{(\Alert[fuchsia]{$R_{\text{stim}}=B_{\text{stim}} N_{2} Y(\Delta E)$})}
\end{itemize}
\end{minipage}

\xdef\loc{east}
\begin{tikzpicture}[remember picture,overlay] % anchor=south
    \node (A) [yshift=0cm,xshift=\figXshift, anchor=\loc] at (current page.\loc) {\includegraphics[page=1,height=8cm,trim={0 0 13cm 0},clip]{Figures_booklet/SOM_9_Statistical_mechanics_figures/SOM_Figure_22.png}}; %scale=\figscale. % 8cm max hei
\end{tikzpicture}


\endofslide 
%\endofsection
\end{frame}
%%%%%%%%%%%%%%%

%%%%%%%%%%%%%%%
\begin{frame}
\frametitle{\texorpdfstring{\culine[\sectiontitleulinecolor]{\sectiontitle}}{\sectiontitle} }
\framesubtitle{Einsteinin kertoimet}
\appear{}

\begin{itemize}
	\item Kertoimet $A_{\text {spon}}$\appear{, $B_{abs}$} \appear{ja $B_{\text {stim}}$} \appear{tunnetaan nimellä \CDAlert[red]{Einsteinin kertoimet} }
	
	\item Tasapainossa näillä prosesseille pätee:
\[
\begin{aligned}
& & R_{\text {spon}} \appear{+R_{\text {stim}}}  &\equal \appear{R_{abs}} \\
 &\appear{\Rightarrow}& \qquad \appear{A_{\text {spon}} N_{2}} \appear{+B_{\text {stim}} N_{2} Y(\Delta E)} & \equal \appear{B_{abs} N_{1} Y(\Delta E)} \qquad \quad 
\end{aligned}
\]

\item Mikä voidaan ratkaista $Y(\Delta E)$ suhteen\appear{, ja käyttämällä $\frac{N_{1}}{\appear{N_{2}}}  \equal \frac{e^{-\Frac{E_{1}}{k_{B} T}}}{\appear{e^{-\Frac{E_{2}}{k_{B} T}}}}$}\appear{, saadaan:}
\[
\begin{aligned}
&\Rightarrow& 
\qquad \appear{Y(\Delta E)} & \equal \frac{A_{\text {spon}} \appear{/ B_{abs}}}{\frac{N_{1}}{N_{2}}  \minus \frac{B_{\text {stim}}}{B_{abs}}}
 \equal 
 \frac{A_{\text {spon}} / B_{abs}}{\frac{e^{-E_{1} / k_{B} T}}{e^{-E_{2} / k_{B} T}}  \minus \frac{B_{\text{stim}}}{B_{abs}}} \\
 \vspace{6mm}
&\Rightarrow&
 \appear{Y(\Delta E)} & \equal \frac{A_{\text {spon}} / B_{abs}}{\appear{e^{\Delta E / k_{B} T}}  \minus \frac{B_{\text {stim}}}{B_{abs}}}
  \equal 
  \frac{A_{\text {spon}} / B_{\text {stim}}}{\Frac{B_{abs}}{B_{\text{stim}}} e^{\Delta E / k_{B} T}-1}  \qquad \quad 
\end{aligned}
\]
\end{itemize}

\endofslide 
%\endofsection
\end{frame}
%%%%%%%%%%%%%%%

%%%%%%%%%%%%%%%
\begin{frame}
\frametitle{\texorpdfstring{\culine[\sectiontitleulinecolor]{\sectiontitle}}{\sectiontitle} }
\framesubtitle{Einsteinin kertoimet}
\appear{}

\begin{itemize}
	\item Parametri $Y(\Delta E)$ ei tarkalleen ottaen ole atomeiden kanssa vuorovaikuttavien fotonien lukumäärä\appear{, vaan fotonikaasun spektrinen energiatiheys (oikealla energiavälillä)}\appear{, joka noudattaa \CDAlert[red]{Planckin lakia}} 
	
	\item Eli voidaan muodostaa kaksi yhtälöä:
\[
\begin{aligned}
Y(\Delta E) &= \Frac{A_{\text {spon}} / B_{abs}}{\Frac{N_{1}}{N_{2}} - \Frac{B_{\text {stim}}}{B_{abs}}}
 = 
 \Frac{A_{\text{spon}} / B_{abs}}{\Frac{e^{-E_{1} / k_{B} T}}{e^{-E_{2} / k_{B} T}}  - \Frac{B_{\text {stim}}}{B_{abs}}}
  = 
  \Frac{A_{\text {spon}} / B_{abs}}{e^{\Delta E / k_{B} T} - \Frac{B_{\text {stim}}}{B_{\text {abs }}}}
   = 
   \Frac{A_{\text {spon}} / B_{\text {stim}}}{\Frac{B_{abs}}{B_{\text {stim}}} e^{\Delta E / k_{B} T}-1} \\
\appear{d U_{fotoni}(\Delta E)} &\equal 
 \frac{h f^{3}}{e^{\Frac{\Delta E}{k_{B} T}}-1} \frac{8 \pi V}{c^{3}} \appear{d f} \quad \Rightarrow \quad \frac{1}{V} \frac{d U_{fotoni}}{d f}
  \equal \frac{8 \pi}{c^{3}} \frac{h f^{3}}{e^{\Frac{\Delta E}{k_{B} T}}-1}  \equal  \appear{Y(\Delta E)}
\end{aligned}
\]
\item Jolloin \CDAlert[blue]{tasapainossa} saadaan muodostettua ehto:
\[
B_{\text {stim}} \equal \appear{B_{abs}} \qquad \appear{\text{jolloin}} \qquad \appear{\frac{A_{\text {spon}}}{\appear{B_{\text {stim}}}}  \equal \frac{8 \pi h f^{3}}{c^{3}}}
\]
\end{itemize}

\endofslide 
%\endofsection
\end{frame}
%%%%%%%%%%%%%%%

%%%%%%%%%%%%%%%
\begin{frame}
\frametitle{\texorpdfstring{\culine[\sectiontitleulinecolor]{\sectiontitle}}{\sectiontitle} }
\framesubtitle{Populaatioinversio}
\appear{}

\begin{itemize}
	\item Ehto \appear{$B_{\text{stim}}=B_{abs}$} \appear{tarkoittaa että on yhtä todennäköistä että tuleva fotoni absorboituu} \appear{tai stimuloi emission}

\item Tämä ehto ei ole laseroinnin kannalta hyödyllinen\appear{, stimuloidun emission pitäisi dominoida}\appear{, koska haluamme valon vahvistumista}\appear{, jonka \CDAlert[red]{stimuloitu emissio mahdollistaa}}
\end{itemize}

\appear{\xdef\loc{south}
\begin{tikzpicture}[remember picture,overlay] % anchor=south
    \node (A) [yshift=-0.25*\figYshift,xshift=0*\figXshift, anchor=\loc] at (current page.\loc) {\includegraphics[page=11,height=5.3cm]{Figures_booklet/SOM_9_Statistical_mechanics_figures/Tikz_SOM_Statistical_figures.pdf}}; %scale=\figscale. % 8cm max height
\end{tikzpicture}}

\endofslide 
%\endofsection
\end{frame}
%%%%%%%%%%%%%%%

%%%%%%%%%%%%%%%%
%\begin{frame}
%\frametitle{\texorpdfstring{\culine[\sectiontitleulinecolor]{\sectiontitle}}{\sectiontitle} }
%\framesubtitle{Toimintaperiaatteet}
%\appear{}
%
%\begin{itemize}
%	\item Laserin toimintaperiaatteet ovat käytännössä:
%
%\item[] \begin{itemize}[itemsep=2.3mm] \small
%	\item Väliaineen reunoilla on samansuuntaiset peilit
%  \item Optisen akselin suhteen eri suuntaan etenevät fotonit \appear{ eivät juurikaan stimuloi emissiota}
%  \item \Alert[blue]{Akselin suuntaiset fotonit heijastelevat peilien välissä}\appear{, mikä mahdollistaa stimuloidun emission}
%  \item Pituus $L$ \appear{ määrittää seisovan aallon:} \appear{$L=n \lambda / 2$}
%    % 6.5 cm = midway, 10 cm = 3/4 
%\vspace{2.3mm}
%\end{itemize}
%\end{itemize}
%\begin{minipage}{6.25cm}
%\begin{itemize}
%\item[] \begin{itemize}[itemsep=2.3mm]  \small
%  \item Kaviteettiin muodostuu resonanssi
%  \item Akselin suuntainen valo vahvistuu
%  \item Erisuuntainen valo ei
%  \item \Alert[red]{Laserin aallonpituus kapenee suhteessa atomien spontaanin transition viivanleveyteen}
%  \item Energiaa tuodaan kaviteettiin esim. 'optisella \Alert[fuchsia]{pumppauksella}'
%\end{itemize}
%\end{itemize}
%  \end{minipage}
%
%\xdef\loc{south east}
%\begin{tikzpicture}[remember picture,overlay] % anchor=south
%    \node (A) [yshift=-0.25*\figYshift,xshift=\figXshift, anchor=\loc] at (current page.\loc) {\includegraphics[page=1,width=6.7cm]{Figures_booklet/SOM_9_Statistical_mechanics_figures/SOM_Figure_23.png}}; %scale=\figscale. % 8cm max hei
%\end{tikzpicture}
%
%
%\endofslide 
%%\endofsection
%\end{frame}
%%%%%%%%%%%%%%%%

%%%%%%%%%%%%%%%
\begin{frame}
\frametitle{\texorpdfstring{\culine[\sectiontitleulinecolor]{\sectiontitle}}{\sectiontitle} }
\framesubtitle{Populaatioinversio}
\appear{}

\begin{itemize}
	\item Normaalissa lämpötilassa\appear{, tasapainossa}\appear{, lähes kaikki atomit ovat perustilassa} 
	\item Eli $N_{1} \gg N_{2}$\appear{, jolloin prosessien nopeuksille pätee}\appear{ $N_{1} B_{abs} \gg N_{2} B_{\text{stim}}$}  
	\implies{ \Alert[red]{Stimuloitu emissio ei dominoi} }

\item Pumppaamalla (optisesti)\appear{}\appear{ perustilassa olevia atomeita voidaan virittää korkeammille energiatiloille} 
\implies{ Korkeamman energiatason $E_{3}$ populaatio kasvaa }

\item Jos vahvistinaineessa \appear{(usein \Alert[tuniblue]{kiinteä}, \CDAlert[tuniblue]{läpinäkyvä kide})} \appear{on \CDAlert[fuchsia]{metastabiili tila  $E_{2}$}}\appear{, niin useat tilalla $E_{3}$ olevat atomit relaksoituvat tälle metastabiilille tilalle $E_{2}$} 
\implies{ On mahdollista saavuttaa tilanne $N_{1}<N_{2}$ }

\end{itemize}

\endofslide 
%\endofsection
\end{frame}
%%%%%%%%%%%%%%%

%%%%%%%%%%%%%%%
\begin{frame}
\frametitle{\texorpdfstring{\culine[\sectiontitleulinecolor]{\sectiontitle}}{\sectiontitle} }
\framesubtitle{Populaatioinversio}
\appear{}

\begin{itemize}[itemsep=1.8mm]
	\item Ensin saavutetaan pumppaamalla tilanne \appear{jossa $N_{1} \approx \appear{N_{2}} \approx \appear{N_{3}}$} \appear{(keskimmäinen kuva)}
	\item Kun pumppaus lopetetaan\appear{, tilan $E_{3}$ atomit alkavat relaksoitumaan}\appear{ kasvattaen tilan $E_{2}$ populaatiota}\appear{, johtaen tilanteeseen $N_{1}<N_{2}$} \appear{(kuva oikealla)}
	
	\item Tätä ilmiötä\appear{/ehtoa} \appear{kutsutaan \CDAlert[red]{populaatioinversioksi}} 

\end{itemize}
\vspace{2.5mm}
\begin{minipage}{3.8cm}
\begin{itemize}
	 \item Kun fotonikaasu alkaa stimuloimaan metastabiilin tilan $E_{2}$ atomeita, \Alert[red]{sti\-mu\-loidun emission nopeus alkaa kasvamaan}
\end{itemize}
\end{minipage}
\vspace{2.5mm}


\xdef\loc{south east}
\begin{tikzpicture}[remember picture,overlay] % anchor=south
    \node (A) [yshift=-0.25*\figYshift,xshift=\figXshift, anchor=\loc] at (current page.\loc) {\includegraphics[page=1,width=9.75cm]{Figures_booklet/SOM_9_Statistical_mechanics_figures/SOM_Figure_26.png}}; %scale=\figscale. % 8cm max hei
\end{tikzpicture}


%\endofslide 
\endofsection
\end{frame}
%%%%%%%%%%%%%%%
